\chapter{Related Work}
\label{ch:related-work}

% Position the thesis precisely. Each subsection ends with "what we add".

\section{Vision-Based Invoice and Document Information Extraction}
\label{sec:rw-extraction}

% TODO: Berghaus et al. 2025 \cite{berghaus2025} -- cloud models on invoices,
% NO recent open-source document VLMs (the gap). Historical: Donut
% \cite{kim2022donut}, LayoutLMv3 \cite{huang2022layoutlmv3}. IDP benchmarks.

\section{Evaluation of Extraction Quality and Measurement Validity}
\label{sec:rw-evaluation}

% TODO: this is where the thesis's methodological contribution is positioned.
% Prior work on scoring extraction against structured references, on
% representation-tolerant matching, and on pre-registration discipline
% \cite{kerr1998}. What this thesis adds: an explicit, documented separation of
% model error from scoring artefact on a task where the reference is machine
% readable while the input is a rendered page, plus an attribution design that
% bounds the structuring stage's contribution using reference-quality transcripts.

\section{Knowledge-Graph Construction and Graph-Augmented Retrieval}
\label{sec:rw-kg}

% KEEP SHORT and be honest about positioning: these works motivated the DESIGN of
% the unbuilt layers; this thesis does not contribute to them. Do not write "what
% we add" here -- there is nothing to add yet. Point forward to
% Ch.~\ref{ch:limitations-future-work} instead.
% TODO: \cite{cai2025} on text-to-graph extraction fidelity at micro and macro
% level, and \cite{han2025} on graph-augmented versus flat retrieval, which
% together shaped the intended Layer-2/3 design and the metrics that would have
% been used.

\section{Positioning of This Thesis}
\label{sec:rw-positioning}

% TODO: one paragraph synthesising the gap this thesis fills.
